Estas son algunas preguntas que surgen de la práctica, con su respuesta.

\textbf{1. ¿Por qué el rayo central no se desvía?}

Porque llega a la superficie justo en la dirección del radio, que es la normal en ese punto. Su ángulo de incidencia es cero y, por la ley de Snell, su ángulo de refracción también.

\textbf{2. ¿Qué pasaría si el cilindro fuera de un material con mayor índice de refracción?}

Los rayos se doblarían más y el foco quedaría todavía más cerca. Según la ecuación \eqref{eq:cilindro}, con un índice de 2 el foco estaría exactamente sobre la cara de salida, ya que $f$ sería igual a $R$.

\textbf{3. ¿Cómo se podría reducir la aberración esférica?}

Usando solo los rayos cercanos al eje, por ejemplo tapando las rendijas de las orillas, o cambiando la forma de la superficie por una que no sea un círculo, como una parábola o una hipérbola, que sí enfocan todos los rayos en un mismo punto.

\textbf{4. ¿Una gota de agua funciona igual?}

Sí, en principio. Una gota es una esfera y concentra la luz de la misma forma, aunque con el índice del agua, de 1.33, su foco queda un poco más lejos. Este mismo efecto, junto con la dispersión, es el que forma el arcoíris.
